% ch7_c §7.6–7.9 (书页 228–238 / p242–p252)
%--B-TAIL--
%JOIN%配位壳层——即径向分布函数 $P(R)$ 中的近邻壳层（参见式 (2.103)）——在 $S(K)$ 中给出一个峰，越过此峰后 $S(K)$ 便在 1 附近振荡。在这张图上，Fermi 球的直径 $2k_{F}$ 依赖于价数 $Z$，因此 (7.72) 中的积分上限在碱金属中达不到 $S(K)$ 的峰，而当 $Z=2$ 时则落入其后的谷中。

由于没有能够产生 Brillouin 区等等的长程序，N 过程与 U 过程之间的区分便失去意义。尽管如此，几何因子 $K^{3}$ 使被积函数的权重强烈地偏向 $2k_{F}$ 附近的值，于是“芯”赝势 $\mathcal{V}'$ 再一次成为占主导的效应。金属熔化时电阻的增大，可以归结为结构因子在大 $K$ 处的变化——即从仍正比于 $T$ 的电子--声子 U 过程，变为来自 $S(K)$ 峰附近的散射；此峰已达 1 的量级，而且在更高的温度下实际上还可能减小。

\section{载流子迁移率}

半导体中载流子被晶格振动的散射，就其一般原理而言是一个简单得多的问题。因为载流子通常被认为集中在 $\mathbf{k}$ 空间的一个小区域内，即 $\mathcal{E}(\mathbf{k})$ 的某个极小值附近，散射中 $\mathbf{k}$ 矢量的可能改变很小。这使我们能够利用 §6.13 的形变势（deformation potential）。散射概率正比于
\begin{equation*}
|\mathcal{M}_{\mathbf{k}'\mathbf{k}}|^{2}\,\mathcal{N}(\mathcal{E})=\Bigl|\sum_{ij}\mathsf{W}^{0}_{ij}\,\mathcal{E}_{ij}\Bigr|^{2}\mathcal{N}(\mathcal{E}),\tag{7.73}
\end{equation*}
其中 $\mathsf{W}^{0}_{ij}$ 是弹性波中应变某个分量的振幅，而 $\mathcal{E}_{ij}$ 是形变势张量的一个分量。

由于 $\mathbf{k}'-\mathbf{k}$ 很小，我们可以对弹性波诸模采用经典统计：如同 (7.54) 那样，
\begin{equation*}
\overline{|\mathsf{W}^{0}_{ij}|^{2}}\propto kT.\tag{7.74}
\end{equation*}
于是，正比于 (7.73) 之倒数的弛豫时间将按如下方式变化
\begin{equation*}
\frac{1}{\tau(\mathcal{E})}\propto T\mathcal{E}^{\frac{1}{2}}m^{*\frac{3}{2}},\tag{7.75}
\end{equation*}
因为态密度是能量和“有效质量”二者的函数（参见式 (4.33)）。把它代入 (7.38) 和 (7.39)，我们得到
\begin{equation*}
\mu\propto T^{-\frac{3}{2}}m^{*-\frac{5}{2}}.\tag{7.76}
\end{equation*}
小有效质量带来的“报酬”是非常高的。

然而，即使引入矩阵元等等，这个理论也仍然太简单。许多半导体的电子极小值不止一个——用行话说，有很多\emph{能谷}（many valleys）（§6.6）——而借助大波矢声子进行的谷间散射十分重要。关系式 (7.76) 在实践中并不很令人满意。

在半导体中，发射或吸收声子时载流子能量的变化不可忽略。它对低场下载流子迁移率的影响不大，但却是\emph{高场}（high-field）输运现象中的主导机制。例如，设想在一个薄的导电样品两端加上几伏的电动势。在两极之间运动的载流子必须把这些能量作为 Joule 热放出——实际上就是发射若干个声子。在金属中，电子--声子两次散射之间的平均自由程非常短（参见式 (7.62)），以致在这个机制失效之前，样品本身早就被气化了。然而在好的半导体中，相当温和的电场（例如 1000 V/cm）就能使能量转移过程过载，于是电子的能量分布可能不再像 (7.38) 所假设的那样处于晶格温度下的平衡。随之发生的\emph{热电子}（hot electron）现象在理论上相当复杂，因为它细致地依赖于电子--声子矩阵元、声学支与光学支的频谱以及固体的电子能带结构。不过，总的效果是偏离 Ohm 定律，趋向于载流子电流几乎与 $E$ 无关的状态。

这个关于电导的讨论还假定晶体动量 $\mathbf{k}$ 是电子态的一个好量子数。在某些情形下——例如在玻璃、非晶态半导体或补偿杂质带中（§§4.2、5.9、6.7）——最好把载流子描述为\emph{定域}（localized）在材料中的各个位置上。这时，导电只能通过热激活的\emph{跳跃}（hopping）从一个位置到另一个位置来实现，同时发射或吸收一个合适的声子以容许这一跃迁。

\section{普遍输运系数}

现在假定样品中除了电场之外还存在温度梯度。忽略尺寸和形状效应，Boltzmann 方程 (7.14) 就成为
\begin{equation*}
\left(-\frac{\partial f^{0}}{\partial\mathcal{E}}\right)\mathbf{v}_{\mathbf{k}}\cdot\left\{\frac{\mathcal{E}(\mathbf{k})-\zeta}{T}(-\nabla T)+e\left(\mathbf{E}-\frac{1}{e}\nabla\zeta\right)\right\}=-\left.\frac{\partial f_{\mathbf{k}}}{\partial t}\right]_{\text{scatt.}}\tag{7.77}
\end{equation*}
我们所做的一切，只是把 (7.16) 中的 $\mathbf{E}$ 换成一个更为复杂的、依赖于 $\mathbf{k}$ 的矢量函数。如果弛豫时间得以存在的条件仍然成立——正如 §7.3 中讨论过的——那么我们就知道解的形式：
\begin{equation*}
f_{\mathbf{k}}-f^{0}_{\mathbf{k}}=\left(-\frac{\partial f^{0}}{\partial\mathcal{E}}\right)\tau\mathbf{v}_{\mathbf{k}}\cdot\left[e\left(\mathbf{E}-\frac{1}{e}\nabla\zeta\right)+\frac{\mathcal{E}(\mathbf{k})-\zeta}{T}(-\nabla T)\right].\tag{7.78}
\end{equation*}
我们可以把它代入一个像 (7.20) 那样的积分中来计算电流，
\begin{align*}
\mathbf{J}&=\frac{1}{4\pi^{3}}\,\frac{e^{2}\tau}{\hbar}\iint\mathbf{v}_{\mathbf{k}}\mathbf{v}_{\mathbf{k}}\left(-\frac{\partial f^{0}}{\partial\mathcal{E}}\right)\frac{\mathrm{d}S}{v_{\mathbf{k}}}\,\mathrm{d}\mathcal{E}\cdot\left(\mathbf{E}-\frac{1}{e}\nabla\zeta\right)\\
&\quad+\frac{1}{4\pi^{3}}\,\frac{e^{2}\tau}{\hbar}\iint\mathbf{v}_{\mathbf{k}}\mathbf{v}_{\mathbf{k}}\left(\frac{\mathcal{E}-\zeta}{T}\right)\left(-\frac{\partial f^{0}}{\partial\mathcal{E}}\right)\frac{\mathrm{d}S}{v_{\mathbf{k}}}\,\mathrm{d}\mathcal{E}\cdot(-\nabla T).\tag{7.79}
\end{align*}
这表明，温度梯度 $\nabla T$ 即使单独作用也会产生电流——这正是存在温差电效应（thermo-electric effect）的证据。

但是，温度梯度最重要的效应是产生\emph{热}流。我们也许会把它看作一种能量通量，
\begin{equation*}
2\int f_{\mathbf{k}}\,\mathcal{E}(\mathbf{k})\,\mathbf{v}_{\mathbf{k}}\,\mathrm{d}\mathbf{k}.
\end{equation*}
（译注：按式 (7.78) 代入 (7.20)，第二项似应为 $e\tau/\hbar$；原书印作 $e^{2}\tau$。）
但这并不正确。“热”是“内能”减去“自由能”。电子的自由能是 $\zeta$，即它的化学势。因此，处于态 $\mathbf{k}$ 的电子所携带的热量是 $\mathcal{E}(\mathbf{k})-\zeta$。更严格的分析证实，热通量的总量（每单位体积）必须由下式计算：
\begin{equation*}
\mathbf{U}=2\int f_{\mathbf{k}}\{\mathcal{E}(\mathbf{k})-\zeta\}\,\mathbf{v}_{\mathbf{k}}\,\mathrm{d}\mathbf{k}.\tag{7.80}
\end{equation*}

在写下 $\mathbf{U}$ 的表达式之前，我们还可以注意到，(7.78) 中含有一个 $\nabla\zeta$ 项——这一项来自电子气温度的改变所引起的化学势变化（参见式 (4.23)）。但在电学性质的普通测量中，任何这类 $\zeta$ 的梯度都应当计入电路中所测得的诸电动势之内。我们将径直略去这一项，假定它的任何效应都已包含在 $\mathbf{E}$——即“观测到的”电场——之中。

由 (7.78)、(7.79) 和 (7.80)，我们可以构造出如下的普遍输运方程
\begin{align*}
\mathbf{J}&=e^{2}\mathsf{K}_{0}\cdot\mathbf{E}+\frac{e}{T}\,\mathsf{K}_{1}\cdot(-\nabla T),\tag{7.81a}\\
\mathbf{U}&=e\,\mathsf{K}_{1}\cdot\mathbf{E}+\frac{1}{T}\,\mathsf{K}_{2}\cdot(-\nabla T),\tag{7.81b}
\end{align*}
其中各系数（其实是张量，不过这只是为了抽象的普遍性而已）定义为
\begin{equation*}
\mathsf{K}_{n}\equiv\frac{1}{4\pi^{3}}\,\frac{\tau}{\hbar}\iint\mathbf{v}\mathbf{v}\,(\mathcal{E}-\zeta)^{n}\left(-\frac{\partial f^{0}}{\partial\mathcal{E}}\right)\frac{\mathrm{d}S}{v}\,\mathrm{d}\mathcal{E}.\tag{7.82}
\end{equation*}
这些系数可以借助我们关于 Fermi 函数对能量积分的普遍定理 (4.21) 来求值：
\begin{equation*}
\int\Phi(\mathcal{E})\left(-\frac{\partial f^{0}}{\partial\mathcal{E}}\right)\mathrm{d}\mathcal{E}=\Phi(\zeta)+\tfrac{1}{6}\pi^{2}(kT)^{2}\left[\frac{\partial^{2}\Phi(\mathcal{E})}{\partial\mathcal{E}^{2}}\right]_{\mathcal{E}=\zeta}+\cdots,\tag{7.83}
\end{equation*}
这里 $\Phi(\mathcal{E})$ 本身是能量面 $\mathcal{E}(\mathbf{k})=\mathcal{E}$ 上的一个积分。但这并不带来特殊的困难。

正如 (7.21) 中那样，我们有
\begin{equation*}
\mathsf{K}_{0}=\frac{1}{4\pi^{3}}\,\frac{\tau}{\hbar}\int\mathbf{v}\mathbf{v}\,\frac{\mathrm{d}S_{F}}{v}.\tag{7.84}
\end{equation*}
现在设想 $\mathsf{K}_{0}(\mathcal{E})$ 是 (7.84) 在某个能量为 $\mathcal{E}$（而非 费米能级处）的面上的取值。对于 $\mathsf{K}_{2}$，我们必须考虑 $\Phi(\mathcal{E})=(\mathcal{E}-\zeta)^{2}\mathsf{K}_{0}(\mathcal{E})$，它在 $\mathcal{E}=\zeta$ 处为零，于是在 (7.83) 的第二项中只剩下
\begin{equation*}
\mathsf{K}_{2}=\tfrac{1}{3}\pi^{2}(kT)^{2}\mathsf{K}_{0}(\zeta)\tag{7.85}
\end{equation*}
而 $\mathsf{K}_{1}$ 的表达式则化为
\begin{equation*}
\mathsf{K}_{1}=\tfrac{1}{3}\pi^{2}(kT)^{2}\left[\frac{\partial}{\partial\mathcal{E}}\,\mathsf{K}_{0}(\mathcal{E})\right]_{\mathcal{E}=\zeta},\tag{7.86}
\end{equation*}
这在原理上要更复杂一些。

(7.81a) 中 $\mathbf{E}$ 的系数当然就是电导率——即我们在等温条件（$\nabla T=0$）下应当测得的电流。我们只是重新确认了 (7.23)；我们写下
\begin{equation*}
\boldsymbol{\sigma}=e^{2}\mathsf{K}_{0},\tag{7.87}
\end{equation*}
而当张量由于具有立方对称性而约化为标量时，还有与之等价的关系式。我们实际上就是利用 (7.87) 用已知的电导率来表示 $\mathsf{K}_{0}$。

\section{热导率}

(7.81b) 是热流方程，其中温度梯度的系数就是我们应当称为电子气的热导率（thermal conductivity）的量。

实际上，这并不完全正确。我们做实验时并不会保持 $\mathbf{E}=\mathbf{0}$。更方便的做法是把样品接成开路，使 $\mathbf{J}=\mathbf{0}$。由 (7.81a)，这意味着沿导线将建立起一个电场
\begin{equation*}
\mathbf{E}=\frac{1}{e^{2}}\,\mathsf{K}_{0}^{-1}\mathsf{K}_{1}\,\frac{e}{T}\cdot\nabla T\tag{7.88}
\end{equation*}
把它代入 (7.81b)，我们得到
\begin{align*}
\mathbf{U}&=\frac{1}{T}\,\mathsf{K}_{1}\mathsf{K}_{0}^{-1}\mathsf{K}_{1}\cdot\nabla T-\frac{1}{T}\,\mathsf{K}_{2}\cdot\nabla T\\
&=\frac{1}{T}\,(\mathsf{K}_{2}-\mathsf{K}_{1}\mathsf{K}_{0}^{-1}\mathsf{K}_{1})\cdot(-\nabla T),\tag{7.89}
\end{align*}
我们应当把它与
\begin{equation*}
\mathbf{U}=\boldsymbol{\kappa}\cdot(-\nabla T)\tag{7.90}
\end{equation*}
中的热导率 $\kappa$ 等同起来。但对金属而言这一修正可以忽略，于是
\begin{equation*}
\kappa=\frac{1}{T}\,\mathsf{K}_{2}.\tag{7.91}
\end{equation*}
而根据 (7.85) 和 (7.87)，此式可以写成
\begin{equation*}
\kappa=\frac{\pi^{2}}{3}\,\frac{k^{2}}{e^{2}}\,T\sigma,\tag{7.92}
\end{equation*}
这个关系称为 Wiedemann--Franz 定律。它是一个老结果，而且非常容易理解。在电传导中，每个电子携带电荷 $e$，受力 $e\mathbf{E}$ 的作用；单位场强下的电流正比于 $e^{2}$。在热传导中，每个电子携带其热能 $kT$，受“热力”$k\nabla T$ 的作用；单位温度梯度下的热流正比于 $k^{2}T$。这两个输运系数之比必须具有 $k^{2}T/e^{2}$ 的量级；因子 $\frac{1}{3}\pi^{2}$ 来自如下事实：我们处理的只是 费米面上服从 费米统计的电子。

容易证明，这个定律是非常普遍的。倘若 (7.40) 成立——也就是说，倘若散射的效应只用一个在 费米面上变化的矢量平均自由程就能定义——这一定律便会成立。但它要求散射是\emph{弹性}的。

为了证明这一点，让我们考察热传导条件下实际分布函数的形式。由 (7.78)，
\begin{equation*}
f_{\mathbf{k}}-f^{0}_{\mathbf{k}}=\left(-\frac{\partial f^{0}}{\partial\mathcal{E}}\right)\left(\frac{\mathcal{E}-\zeta}{T}\right)\tau\,\mathbf{v}_{\mathbf{k}}\cdot(-\nabla T).\tag{7.93}
\end{equation*}
让我们看一看 费米分布的一个截面。如图 126 所示，在 $f^{0}$ 上加上一个像 (7.93) 那样的表达式，其效果是使 $\mathbf{v}_{\mathbf{k}}\cdot(-\nabla T)$ 为正的一侧的分布展宽，而在另一侧则使分布变窄。

我们可以把它解析地表示出来：
\begin{align*}
f_{\mathbf{k}}&=f^{0}_{\mathbf{k}}-(\tau\,\mathbf{v}_{\mathbf{k}}\cdot\nabla T)\,\frac{\partial f^{0}}{\partial T}\\
&=f^{0}(T-\tau\,\mathbf{v}_{\mathbf{k}}\cdot\nabla T).\tag{7.94}
\end{align*}
沿 $\nabla T$ 为负的方向运动的电子，即沿温度梯度下行方向运动的电子，要“热”出这样一个数量：
\begin{equation*}
\delta T=-\tau\,\mathbf{v}_{\mathbf{k}}\cdot\nabla T.\tag{7.95}
\end{equation*}
沿相反方向运动的电子，则比电子气的平均温度“冷”。

%FIG{126}: {热传导中的 费米分布。}

的确，这正是我们按分子运动论的推理所应得出的结论。从方向 $\mathbf{v}$ 过来、到达温度为 $T$ 的区域的电子，已经走过了一段距离 $\tau v$。它们离开的那个区域——即它们上一次经受使其热化的碰撞的区域——将处于温度 $T+\delta T$。因此，这些电子要“热”出一个量 $\delta T$。

于是，在热传导实验中没有净的电子通量——没有净的电荷通量。之所以有热流，是因为有“热”电子向一个方向行进，同时有同等数目的“冷”电子向相反方向行进。图 127(a) 中粗略地画出了这一情形。右边有一些被激发到 费米能级以上的电子；左边则有一些凝聚入 费米面之内的电子。

现在考虑这个分布可以通过哪些散射方式而弛豫。我们可以把一个热电子沿 Fermi 球一直散射过去，使其方向反转。我们称之为\emph{水平过程}（horizontal process）。这类散射同样也是降低电导率的有效过程，如图 127(b) 所示。这类过程是弹性的——对它们而言 Wiedemann--Franz 定律成立。

但也存在\emph{竖直}跃迁：一个“热”电子失去它的全部多余能量，落到 费米能级以下。这类过程对电导率的影响很小——但在热传导中，它在削减热流方面与大角散射同等重要。

%FIG{127}: {电子分布与散射过程：(a) 热传导中；(b) 电传导中。}

按定义，竖直跃迁是\emph{非弹性}的。杂质散射不会发生这种跃迁，由杂质散射引起的热阻服从 Wiedemann--Franz 定律。换言之，由于 $\rho$ 为常数，$\kappa\propto T$。在高温下的声子散射中，竖直跃迁也不十分重要，因为那时电子能够获得或失去的最大能量就是一个声子的最大能量，即 $k\Theta$。它小于 费米分布相对于 费米能级的展宽 $kT$：结果表明，在高温极限下，对“理想”电阻而言 Wiedemann--Franz 定律同样成立。由于 $\rho_{i}$ 正比于 $T$，我们发现 $\kappa_{i}$ 应当趋于一个常数。

但在低温下，电子平均获得或失去的能量就是该温度下声子的平均能量，即它必须具有 $kT$ 的量级。这恰好足以使电子穿过热层——恰好足以把一个“热”电子变成“冷”的。

对这种情形推导一个公式，显然需要仔细研究产生或吸收声子的各个具体过程。但是，如果我们假定现在所有散射过程在削减热流方面都同等有效，就可以大致看出会发生什么。也就是说，在像 (7.46) 那样计算“弛豫时间”时，我们现在略去因子 $(1-\cos\theta)$，正是这个因子顾及了电子速度沿场方向分量的变化。

但这个因子 $(1-\cos\theta)$ 在低温下曾在积分 (7.70) 中贡献了 $2\sin^{2}\frac{1}{2}\theta$。于是，在我们的“截止”近似中应当去掉一个因子 $(T/\Theta)^{2}$。在像 (7.67) 那样的计算中，“热导率弛豫时间”将得出为
\begin{equation*}
\frac{1}{\tau}\propto\frac{T^{3}}{M\Theta^{4}}\tag{7.96}
\end{equation*}
而不是电导率中所得到的 $T^{5}$ 定律。把它代入 (7.91)、(7.85) 和 (7.84)（即在把热导率与电导率相比较时总是要计入的、正比于 $T$ 的 Wiedemann--Franz 因子），我们得到
\begin{equation*}
\kappa_{i}\propto\frac{M\Theta^{4}}{T^{2}}.\tag{7.97}
\end{equation*}
这在实际中被观察到——在 Debye 温度以下，电子--声子过程给出如下形式的 Lorenz 比（Lorenz ratio）：
\begin{equation*}
\frac{\kappa_{i}}{\sigma_{i}T}\propto\left(\frac{T}{\Theta}\right)^{2}\tag{7.98}
\end{equation*}
而不是 Wiedemann--Franz 定律 (7.92) 所要求的常数值 $(\tfrac{1}{3}\pi^{2})(k/e)^{2}$。但在极低温度下，剩余电阻（§7.4）来自杂质的弹性散射，这个比值便回到 Lorenz 比 (7.92)。

当然，还必须考虑 U 过程对热阻率的贡献，但这又是一个复杂的计算问题。

\section{温差电效应}

普遍关系 (7.81) 表明，电流与热流之间存在交互效应。例如，设想在一个开路的样品中建立温度梯度。在 (7.81a) 中令 $\mathbf{J}=\mathbf{0}$（这里悄悄换成标量系数），我们得到
\begin{align*}
\mathbf{E}&=\frac{1}{eT}\,(\mathsf{K}_{0}^{-1}\mathsf{K}_{1})\,\nabla T\\
&=Q\nabla T;\tag{7.99}
\end{align*}
这时将观测到一个电动势。

为了演示并测量这个效应，必须用两种金属 $A$ 和 $B$ 组成一个闭合回路，两个接头处于不同的温度 $T_{1}$ 和 $T_{2}$，并在某个中间点（温度 $T_{0}$）接入一个伏特计。回路中的电动势定义为 $\mathbf{E}$ 沿导线全长的积分
\begin{align*}
\phi&=\int_{0}^{1}E_{B}\,\mathrm{d}x+\int_{1}^{2}E_{A}\,\mathrm{d}x+\int_{2}^{0}E_{B}\,\mathrm{d}x\\
&=\int_{2}^{1}Q_{B}\,\frac{\partial T}{\partial x}\,\mathrm{d}x+\int_{1}^{2}Q_{A}\,\frac{\partial T}{\partial x}\,\mathrm{d}x\\
&=\int_{T_{1}}^{T_{2}}(Q_{A}-Q_{B})\,\mathrm{d}T.\tag{7.100}
\end{align*}
回路中产生的电压看起来像是两个接头温度之差、以及两种金属的\emph{绝对温差电功率}（absolute thermo-electric power）$Q$ 之差的函数。这就是著名的 Seebeck 效应。

%FIG{128}: {Seebeck 效应。}

但是还有另一种现象，它与 (7.81) 中出现的量 $\mathsf{K}_{1}$ 相联系。设想我们使一个如图 129 那样的回路保持 $\nabla T=\mathbf{0}$。由 (7.81a)，我们有
\begin{equation*}
\mathbf{U}=e\,\mathsf{K}_{1}\,\mathbf{E},\qquad\mathbf{J}=e^{2}\mathsf{K}_{0}\,\mathbf{E},\tag{7.101}
\end{equation*}
于是
\begin{align*}
\mathbf{U}&=\frac{1}{e}\,\mathsf{K}_{0}^{-1}\mathsf{K}_{1}\,\mathbf{J}\\
&=\Pi\,\mathbf{J}.\tag{7.102}
\end{align*}
现在用电池在回路中驱动一个电流 $J$。在支路 $A$ 中将有一个热流 $\Pi_{A}J$；在支路 $B$ 中则将有另一个热流 $\Pi_{B}J$。在接头处，热的平衡必须重新建立；热通量 $(\Pi_{A}-\Pi_{B})J$ 将在一个接头处放出，而在另一个接头处吸收。一个接头将变热，另一个变冷。这就是 Peltier 效应。

正如从 (7.99) 和 (7.102) 所看到的，\emph{Peltier 系数}与绝对温差电功率有以下关系：
\begin{equation*}
\Pi=QT.\tag{7.103}
\end{equation*}
这并非偶然。它是温差电学的 \emph{Kelvin 关系}（Kelvin relations）之一，也是不可逆过程热力学的 Onsager 关系的一个特例。实质上，方程 (7.81) 若用 $\mathbf{J}$ 和 $\mathbf{U}/T$ 来表达，就必须具有一个对称的矩阵——这解释了为什么 $\mathsf{K}_{1}$ 作为系数出现了两次。还有另一种
%FIG{129}: {Peltier 效应。}
温差电效应——\emph{Thomson 效应}（Thomson effect），它在同时载有热流和电流的回路中被观察到——不过这个效应的系数已经包含在普遍关系 (7.81) 之中。

至于温差电功率的实际数值，我们援引 (7.86) 和 (7.87)：
\begin{equation*}
Q=\frac{\pi^{2}}{3}\,\frac{k}{e}\,kT\left[\frac{\partial\ln\sigma(\mathcal{E})}{\partial\mathcal{E}}\right]_{\mathcal{E}=\zeta},\tag{7.104}
\end{equation*}
其中 $\sigma(\mathcal{E})$ 意指“在一个假想的、费米能级位于 $\mathcal{E}$ 处的金属中电导率所取的值”。把这个公式应用于各种情形下电导率的计算结果，在原则上足够容易（尽管在实践中往往复杂而不成功）。

例如，设想我们写出
\begin{equation*}
\sigma(\mathcal{E})=n(\mathcal{E})\,e\,\mu(\mathcal{E})\tag{7.105}
\end{equation*}
即以一个随能量变化的载流子密度和一个有效载流子迁移率来表述。于是我们有
\begin{equation*}
Q=\frac{\pi^{2}}{3}\,\frac{k}{e}\,kT\left[\frac{\mathcal{N}(\mathcal{E})}{n}+\frac{\partial\ln\mu(\mathcal{E})}{\partial\mathcal{E}}\right]_{\mathcal{E}=\zeta}.\tag{7.106}
\end{equation*}
第一项很容易理解。正如 §4.7 中所见，“每电子比热”是
\begin{equation*}
C_{\text{el.}}=\frac{\pi^{2}}{3}\,k^{2}T\,\frac{\mathcal{N}(\mathcal{E}_{F})}{n}.\tag{7.107}
\end{equation*}
因此，Peltier 系数 $QT$ 就是“每电子热”与电子电荷之比——正如 (7.100) 所提示的那样。

之所以出现迁移率的导数，是因为我们必须顾及电子电流按能量分布的方式。于是，如果 $\mu(\mathcal{E})$ 随能量增大，那么电流中有很高的比例是由能量较高的电子承载的，而这些电子按比例携带着较大的热流。

当然，对于金属，(7.105) 并不合理；应当像 (7.25) 那样用 费米面面积和平均自由程来表示电导率，或者像 (7.34) 那样用态密度、Fermi 电流和弛豫时间来表示。对于自由电子系统，两种作法的结果相同；但要正确解释金属或合金中的温差电现象，就必须小心选择模型，使得上述各种因素随能量的变化都能得到正确的评估。磁性杂质所产生的巨大温差电功率（§10.6）就是这类效应的一个极端例子。值得注意，电荷 $e$ 按惯例取负值，因此 $Q$ 应当为负——尤其是因为 $\mu(\mathcal{E})$ 倾向于随 $\mathcal{E}$ 而增大；较快的电子不那么容易被散射。

在半导体中，这类公式不能令人满意。我们必须回到各系数的定义 (7.82)，并且不能再像 (7.83) 那样利用 Fermi 函数的性质。特别地，$\mathsf{K}_{1}$ 的积分具有如下形式
\begin{equation*}
\int(\mathcal{E}-\zeta)\,f_{\mathbf{k}}\,\mathbf{v}_{\mathbf{k}}\,\mathrm{d}\mathbf{k},\tag{7.108}
\end{equation*}
其中将含有对应于 §4.6 中 $\zeta_{e}$ 的“热”——即载流子的能量，不是从最近的带边量起，而是从真正的 费米能级量起。于是，温差电功率中有一个形如下式的贡献
\begin{equation*}
Q_{e}\sim\frac{1}{T}\,\frac{\zeta_{e}}{e}\tag{7.109}
\end{equation*}
此外还有一项通常较小的、来自 $\mu(\mathcal{E})$ 行为的贡献。

如果我们的半导体主要含有“空穴”，那么就应当
